\documentclass[12pt]{amsart}

\usepackage{amssymb}
\usepackage{amsthm}
\usepackage{amsmath}
\usepackage{amsxtra}
\usepackage{mathrsfs}
\usepackage{enumerate}
\usepackage{colonequals}
\usepackage{graphicx}

\setlength{\hfuzz}{4pt}

\newenvironment{problab}[1]
{\noindent\textbf{Problem #1}.}
{\vskip 6pt}

\newenvironment{exolab}[1]
{\noindent\textbf{Exercise #1}.}
{\vskip 6pt}


\theoremstyle{remark}
\newtheorem*{soln}{Solution}

\newcommand{\C}{\mathbb C}
\newcommand{\D}{\mathfrak D}
\renewcommand{\H}{\mathfrak H}
\newcommand{\F}{\mathbb F}
\renewcommand{\P}{\mathbb P}
\newcommand{\PP}{\mathbb P}
\newcommand{\Q}{\mathbb Q}
\newcommand{\R}{\mathbb R}
\renewcommand{\SS}{\mathbb S}
\newcommand{\Z}{\mathbb Z}

\newcommand{\wh}{\widehat}

\renewcommand{\Im}{\operatorname{Im}}
\renewcommand{\Re}{\operatorname{Re}}

\newcommand{\eps}{\epsilon}

\usepackage{fullpage}

\DeclareMathOperator{\Aut}{Aut}
\DeclareMathOperator{\Hom}{Hom}
\DeclareMathOperator{\ord}{ord}
\DeclareMathOperator{\Stab}{Stab}

\usepackage{mathpazo}

\begin{document}

\title{Belyi Maps and Dessins d'Enfants \\ Homework \#6}
\date{\today}

%\thanks{\emph{Date:} Monday, 11 February 2013}

\maketitle

\begin{exolab}{6.1}
  Let $T$ be a hyperbolic triangle with one vertex at infinity, and the other two vertices on the unit circle. Denote the angles of $T$ at these latter two vertices by $\alpha$ and $\beta$, as shown below. Compute the hyperbolic area of $T$. (Recall that the hyperbolic area is given by
  \begin{equation*}
    a(T) = \int_T \frac{dx dy}{y^2} \, .
  \end{equation*}
  Use polar coordinates to compute this integral.)
  \vskip 5pt
  \begin{center}
    \includegraphics[width=0.6\linewidth]{triangle-infinity.png}
  \end{center}
\end{exolab}

\begin{exolab}{6.2}
  Let $G$ be a transitive permutation group on a finite set $A$. Recall the following definitions.
  \begin{itemize}
    \item 
      A \textsf{block} of $G$ is a nonempty subset $B$ of $A$ such that for all $\sigma \in G$, either $\sigma(B) = B$ or $\sigma(B) \cap B = \varnothing$, where $\sigma(B) = \{\sigma(b) : b \in B\}$.
    \item
      $G$ is \textsf{primitive} if the only blocks of $G$ are are the trivial ones: the sets of size $1$, and $A$ itself.
  \end{itemize}

  \begin{enumerate}[(a)]
    \item 
      If $B$ is a block and $a \in B$, show that the set
      \begin{equation*}
        \Stab_G(B) = \{\sigma \in G \mid \sigma(B) = B\}
      \end{equation*}
      is a subgroup of $G$ containing $\Stab_G(a)$.
    \item
      If $B$ is a block and $\sigma_1(B), \sigma_2(B), \ldots, \sigma_n(B)$ are the distinct images of $B$ under elements of $G$, show that these form a partition of $A$.
    \item
      Prove that $G$ is primitive if and only if for each $a \in A$, the only subgroups of $G$ containing $\Stab_G(a)$ are $\Stab_G(a)$ and $G$ itself, i.e., $\Stab_G(a)$ is a \emph{maximal} subgroup of $G$.
  \end{enumerate}

\end{exolab}

\end{document}
